\section{Covariance concavity}
\subsection{The simplex dual and its strata}
First let $\nu=\sum_{i=0}^d p_i\delta_{y_i}$, where the sites are affinely
independent and $p_i>0$. For $C>0$,
\begin{equation}\label{eq:dual}
 \mathcal F_\nu(C)=\min_h
 \left\{\E\max_i(y_i\cdot X+h_i)-\sum_i p_i h_i\right\},
 \qquad X\sim N(0,C).
\end{equation}
Indeed the maximum bounds every coupling. The objective is at least
$\max h-\sum p_i h_i\geq(\min p_i)(\max h-\min h)$, so a minimum exists
modulo common offsets. Ties have Gaussian measure zero. Differentiation in
$h_i$ makes the winning masses equal to $p_i$, and that winning assignment
attains the bound.

Every offset class has a unique representation $h_i=c-y_i\cdot a$.
The cells are therefore the target simplex's normal fan translated by $a$.
For each nonempty vertex set $S$, let $\Gamma_S$ be the locus where all
labels in $S$ tie for the maximum. It is a simplicial cone in an affine
space of dimension $d+1-|S|$, based at $a$, with boundary facets
$\Gamma_{S\cup\{k\}}$, $k\notin S$. At $C=I$ define
\[
 \mu_S=\int_{\Gamma_S}\gamma_d(x)\dd\mathcal H^{d+1-|S|}(x),\qquad
 m_S=\frac1{\mu_S}\int_{\Gamma_S}x\gamma_d(x)\dd\mathcal H^{d+1-|S|}(x),
\]
and let $P_S$ be the orthogonal projection onto its tangent space.
All these masses are positive and all polynomial moments are finite.
For the apex, $\mu_{\{0,\ldots,d\}}=\gamma_d(a)$ and its mean is $a$.

The offset objective is the Gaussian convolution of the fixed support
function $\max_i y_i\cdot x$, evaluated at $-a$, plus
$(\sum_i p_i y_i)\cdot a$. Linear growth of the support function gives smooth
dependence on $a,C>0$ by differentiation of the Gaussian density. The offset
Hessian is the complete graph Laplacian with positive facet weights
$\mu_{\{i,j\}}/|y_i-y_j|$. It is invertible modulo constants. The implicit
function theorem thus gives smooth minimizing offsets and a well-defined
optimized Hessian.

\subsection{Intrinsic Gaussian integration}
For $i\in S$ and $k\notin S$, the inward normal within the stratum is
\[
 n_{S,k}=\frac{P_S(y_i-y_k)}{|P_S(y_i-y_k)|}.
\]
Intrinsic integration by parts yields, for positive-dimensional strata,
\begin{equation}\label{eq:stratum-mean}
 \mu_S P_Sm_S=\sum_{k\notin S}\mu_{S\cup\{k\}}n_{S,k}.
\end{equation}
For an edge stratum $e=\{i,j\}$ a second integration gives
\begin{equation}\label{eq:stratum-cov}
 \mu_e(P_e-\Cov_e X)
 =\sum_{k\notin e}\mu_{e\cup\{k\}}
       n_{e,k}(m_e-m_{e\cup\{k\}})^T.
\end{equation}
To verify the signs, the intrinsic Gaussian derivative is $-P_Sx\gamma_d$,
whereas the outward boundary normal is $-n_{S,k}$. In the second identity,
integrate the derivative of $((x-m_e)\cdot u)\gamma_d$ in a tangent direction
$v$. The two negative boundary signs produce the displayed product.
Both means lie in the same affine facet, so their difference is tangent.
Truncation by intrinsic balls produces only polynomial Gaussian tail terms
on the artificial sphere, which vanish as the radius tends to infinity.
Lower-dimensional intersections have zero boundary measure.

We use the classical convex-support Brascamp--Lieb consequence
\begin{equation}\label{eq:BL}
 \Cov_e X\leq P_e.
\end{equation}
For a direct applicable formulation, Bobkov--Ledoux
\cite[Proposition 2.1]{BL2000} allows a smooth strictly convex potential on
an open convex support. In orthonormal intrinsic coordinates the relative
interior of the facet is such a support, and its Gaussian potential has
Hessian $I$. Apply that proposition to linear functions; the facet boundary
has zero intrinsic measure. Neither strictness nor invertibility of the
covariance deficit is needed.

\subsection{Facet matrices}
Orient each edge $e=(i,j)$ with $i<j$ and put
\[
 d_e=y_i-y_j,\qquad \ell_e=|d_e|,\qquad w_e=\mu_e/\ell_e.
\]
For every target triangle $g$ let $A_g>0$ be its Euclidean area and set
$D_g=2A_g/\mu_g$. Let $l_k$ be the global affine barycentric coordinate
of vertex $k$ and write $v_k=-\nabla l_k$. Then
\begin{equation}\label{eq:barycentric}
 \sum_kv_k=0,\qquad v_k\cdot(y_i-y_j)=\delta_{kj}-\delta_{ki}.
\end{equation}
For $k\notin e$, $v_k$ is tangent to $\Gamma_e$. The vectors
$P_e(y_i-y_k)$ form a basis of that tangent space, with dual rows $v_k^T$.

Index the columns of $N_e,T_e$ by triangles $g=e\cup\{k\}$ and define
\[
 N_e[:,g]=\frac{\mu_g}{\ell_e}n_{e,k},\qquad
 T_e[:,g]=m_e-m_g,\qquad E_e=w_e(P_e-\Cov_e X).
\]
All matrices and inverses in this paragraph act on the facet tangent space.
The triangle height formula gives
\begin{equation}\label{eq:Ninverse}
 N_e[:,g]=\frac{P_e(y_i-y_k)}{D_g},
 \qquad \operatorname{row}_g(N_e^{-1})=D_gv_k^T.
\end{equation}
Equations~\eqref{eq:stratum-cov}--\eqref{eq:BL} imply
\begin{equation}\label{eq:Blocal}
 E_e=N_eT_e^T=T_eN_e^T\geq0,\quad
 T_e=E_eN_e^{-T},\quad
 B_e=N_e^{-1}E_eN_e^{-T}=N_e^{-1}T_e\geq0.
\end{equation}
There is no inverse of $E_e$ here. Equation~\eqref{eq:stratum-mean} also gives
$w_eP_em_e=N_e\1$, hence the exact entries
\begin{equation}\label{eq:Bentries}
 B_e[g,g']=\frac1{w_e}-D_gv_k\cdot m_{g'},\qquad g=e\cup\{k\}.
\end{equation}

\subsection{The chain factorization}
Let $\partial_2$ and $\partial_3$ be the usual boundary maps from oriented
triangles to edges and from tetrahedra to triangles, on the complete simplex.
All vertices in a simplex are ordered increasingly; for example
$\partial_2[i,j,k]=[j,k]-[i,k]+[i,j]$.
Coning chains to one fixed vertex proves that the graph cycle space is
$\im\partial_2$ and $\ker\partial_2=\im\partial_3$.
In dimension two $\partial_3$ is absent.

Put $W=\diag(w_e)$ and
$A=\partial_2^TW^{-1}\partial_2$. For a triangle coefficient vector
$\lambda$, let $\eta_{e,g}=(\partial_2)_{e,g}\lambda_g$ on triangles
incident to $e$, and define $B$ by
$\lambda^TB\lambda=\sum_e\eta_e^TB_e\eta_e$.
For a target tetrahedron $h$, let $V_h$ be its intrinsic three-dimensional
volume, and set $c_h=\mu_h/(6V_h)$. Also write $D=\diag(D_g)$.

\begin{lemma}[Exact boundary factorization]\label{lem:factor}
\begin{equation}\label{eq:factor}
 A-B=-D\partial_3\diag(c_h)\partial_3^TD.
\end{equation}
\end{lemma}
\begin{proof}
Suppose distinct triangles $g,g'$ share edge $e$, and let $k=g\setminus e$.
Because $k\notin g'$, the vector $v_k$ is tangent to $\Gamma_{g'}$.
Test~\eqref{eq:stratum-mean} there against $v_k$.
By~\eqref{eq:barycentric}, all boundary normals vanish in this pairing except
the one adjoining $k$. Its pairing is the reciprocal target height
$A_{g'}/(3V_h)$, where $h=g\cup g'$. Thus
\[
 v_k\cdot m_{g'}=\frac{\mu_hA_{g'}}{3V_h\mu_{g'}}.
\]
The off-diagonal entry of $A-B$ is therefore
$(\partial_2)_{e,g}(\partial_2)_{e,g'}D_gD_{g'}c_h$.
The identity $\partial_2\partial_3=0$, on the shared edge of this tetrahedron,
gives
\[
 (\partial_2)_{e,g}(\partial_2)_{e,g'}
 =-(\partial_3)_{g,h}(\partial_3)_{g',h}.
\]
This proves the claimed off-diagonal entries. Triangles sharing no edge
interact in neither matrix.

For the diagonal, summing~\eqref{eq:Bentries} over the three edges of $g$
and using~\eqref{eq:barycentric} gives
\[
 (A-B)_{g,g}
 =D_g\sum_{k\in g}v_k\cdot m_g
 =-D_g\sum_{l\notin g}v_l\cdot m_g
 =-D_g^2\sum_{h\supset g}c_h.
\]
The last equality uses the same stratum mean identity and reciprocal height.
If $d=2$, the sum over outside vertices is empty and the difference is zero,
also when the triangle stratum is the apex. These are all entries of
\eqref{eq:factor}. The volumes in this argument are intrinsic to the target
tetrahedra, so there is no factor depending on the ambient dimension.
\end{proof}

\subsection{An admissible gauge and the cycle estimate}
For a fixed edge cycle $j$, choose $\partial_2\lambda=j$.
Its representatives form $\lambda+\im\partial_3$. Minimize
$\lambda^TD\lambda$ on this closed affine space. Since $D>0$, the minimizing
representative exists and satisfies
\begin{equation}\label{eq:gauge}
 \partial_3^TD\lambda=0.
\end{equation}
This uses no injectivity of $\partial_3$. It changes only cycle coordinates.
By~\eqref{eq:factor},
\begin{equation}\label{eq:gauge-equality}
 \lambda^TB\lambda=\lambda^TA\lambda=j^TW^{-1}j.
\end{equation}

Give the cells arbitrary vectors $z_i$ and write
$z_e=z_i-z_j$, $a_e=z_e\cdot m_e$.
Signed edge differences telescope around each triangle, so
\[
 j\cdot a=\sum_e z_e\cdot T_e\eta_e.
\]
By~\eqref{eq:Blocal}, square-root Cauchy--Schwarz in the PSD matrices $E_e$,
followed by Cauchy--Schwarz across edges, gives
\[
 (j\cdot a)^2\leq
 \left(\sum_e z_e^TE_ez_e\right)j^TW^{-1}j.
\]
Singular deficits cause no exception. Weighted least-squares duality on the
edge graph now yields
\begin{equation}\label{eq:cycle-bound}
 \min_b\sum_e w_e(a_e+b_i-b_j)^2
 \leq\sum_e w_e\bigl(|P_ez_e|^2-\Var_e(z_e\cdot X)\bigr).
\end{equation}
For clarity, at the least-squares minimizer the weighted residual is a cycle.
Cauchy--Schwarz gives the dual supremum
$\sup_{j\ne0}(j\cdot a)^2/(j^TW^{-1}j)$, and that residual attains it
unless it is zero. Thus~\eqref{eq:cycle-bound} uses a single global offset
vector $b$, not offsets chosen independently on facets.

\subsection{The exact covariance Hessian}
Take $C(t)=I+tH$ with $H=H^T$. The first two derivatives of its square root
at zero are $H/2$ and $-H^2/4$. Differentiating the affine maximum in
\eqref{eq:dual} and eliminating the offset derivative gives
\begin{equation}\label{eq:cov-Hessian}
 4D^2\mathcal F_\nu(I)[H,H]
 =\min_b\sum_e w_e\E\bigl[(X\cdot Hd_e+b_i-b_j)^2\mid\Gamma_e\bigr]
       -\sum_e w_e|Hd_e|^2.
\end{equation}
Here conditional expectation means the normalized intrinsic facet measure.
To detail the second variation, the score velocity difference across a facet
is $X\cdot Hd_e/2+h_i'-h_j'$. Differentiating the first derivative over moving
cells contributes its square divided by $\ell_e$, integrated over the facet.
Each unoriented facet is counted once. Gaussian integration by parts for the
winning vector gives
$\E[XY(X)^T]=\sum_e w_ed_ed_e^T$, so the interior acceleration is
$-\frac14\sum_e w_e|Hd_e|^2$. Put $b=2h'$.
The $h''$ term vanishes because the current masses equal $p$.
The positive offset Hessian makes the displayed minimum exactly the Schur
complement. Polyhedral truncation and Gaussian tails justify the boundary
differentiation; multiple ties have zero facet measure, and the smoothness
already proved identifies this calculation with the ordinary Hessian.

Apply~\eqref{eq:cycle-bound} with $z_i=Hy_i$. The facet variances cancel,
leaving
\begin{equation}\label{eq:strict-Hessian}
 4D^2\mathcal F_\nu(I)[H,H]
 \leq-\sum_e w_e\frac{(d_e^THd_e)^2}{\ell_e^2}.
\end{equation}
The edge dyads of an affinely independent simplex span the symmetric matrices:
differences $b_i=y_i-y_0$ form a basis, and the quadratic forms at $b_i$
and $b_i-b_j$ recover every bilinear entry in that basis.
Thus the right side is strictly negative for $H\ne0$.
At any $C_0>0$, transform the target sites by $C_0^{1/2}$ and the covariance
direction $K$ by $H=C_0^{-1/2}KC_0^{-1/2}$. This reduces to the calculation
at $I$, without a commutativity assumption. Integration along segments proves
concavity on the positive definite cone. For $d=1$ the value is a positive
constant times the square root of variance, giving the same assertions.

\subsection{Singular covariance, finite lifting, and limits}
For every fixed $\nu\in\cP_2$,
\begin{equation}\label{eq:C-continuity}
 |\mathcal F_\nu(C)-\mathcal F_\nu(C')|
 \leq(\E|Y|^2)^{1/2}\norm{C^{1/2}-(C')^{1/2}}_{\HS}.
\end{equation}
Indeed any coupling of $C^{1/2}G$ with $Y$ lifts to one with standard $G$,
using its Gaussian conditional law even when $C$ is singular. Reuse $G,Y$
for $C'$ and apply Cauchy--Schwarz; then reverse the argument.
Adding $\eps I$ and taking limits extends simplex concavity to the PSD cone.

For an arbitrary finite target $\sum_{i=1}^m p_i\delta_{y_i}$ in $\R^N$,
lift its sites to $(y_i,\eps e_i)$ in $\R^{N+m}$, for any fixed $\eps>0$.
They are affinely independent. A Gaussian with covariance $\diag(C,0)$
has exactly the old cost, so the two optimal values coincide. Subtract a
common lifted site; its expected contribution is zero. Project onto the
$(m-1)$-dimensional span $S$ of lifted differences. The projected source has
covariance $P_S\diag(C,0)P_S$, linear in $C$, and its cost is unchanged.
Every projected coupling lifts by the source conditional law. Simplex PSD
concavity therefore pulls back to the original finite-target functional.
A one-atom target has value zero.

Finally, gluing a coupling of two targets to a Gaussian-target coupling gives
\begin{equation}\label{eq:target-continuity}
 |\mathcal F_\nu(C)-\mathcal F_{\nu'}(C)|
 \leq\sqrt{\tr C}\,W_2(\nu,\nu').
\end{equation}
Finite laws are dense in $\cP_2$, by truncation and quantization. Pass to the
limit in the concavity inequality. The identity
$W_2^2(N(0,C),\nu)=\tr C+\int|y|^2\dd\nu-2\mathcal F_\nu(C)$ proves the
equivalent assertion in Theorem~\ref{thm:covariance}. Neither lifting nor
limits assert strictness for general targets.

\begin{corollary}\label{cor:square}
The function $C\mapsto\mathcal F_\nu(C)^2$ is concave on the PSD cone.
\end{corollary}
\begin{proof}
Independence gives $\mathcal F_\nu\geq0$, and scaling the Gaussian gives
$\mathcal F_\nu(tC)=\sqrt t\,\mathcal F_\nu(C)$.
For $0<a<1$, concavity applied to $C/a$ and $D/(1-a)$ yields
\[
 \mathcal F_\nu(C+D)\geq
 \sqrt a\,\mathcal F_\nu(C)+\sqrt{1-a}\,\mathcal F_\nu(D).
\]
Optimizing $a$, with endpoints obtained by limits, proves superadditivity
of $\mathcal F_\nu^2$. Its degree-one homogeneity then proves concavity.
\end{proof}

